\begin{solapartat}
La força d'atracció gravitatòria de Saturn sobre Mimas serà igual a la massa de Mimas per la seva acceleració centrípeta: \punts{0,1}
\[ G\,\frac{M_S\,M_M}{r_\text{òrbita}^2} = M_M\left(\frac{2\pi}{T}\right)^2 r_\text{òrbita} \quad\punts{0,1} \]
\[ r_\text{òrbita} = \sqrt[3]{\frac{G\,M_S\,T^2}{4\pi^2}} = 1,85\cdot 10^{8}\ \mathrm{m} \quad\punts{0,3} \]
Per tant l'alçada respecte la superfície de Saturn serà: $h = r_\text{òrbita} - R_\mathrm{Saturn} = 1,28\cdot 10^{8}\ \mathrm{m}$ \quad\punts{0,2}

La seva velocitat orbital serà: $v_\mathrm{orbital} = \omega\, r_\text{òrbita} = \frac{2\pi}{T}\, r_\text{òrbita} = 1,43\cdot 10^{4}\ \mathrm{m/s}$ \quad\punts{0,3}
\end{solapartat}

\begin{solapartat}
L'energia mecànica de Mimas serà:
\[ \left.\begin{aligned}
E_m = E_T &= \tfrac{1}{2}\,M_M\,v_\mathrm{orbital}^2 - \tfrac{G\,M_S\,M_M}{r_\text{òrbita}} \quad\punts{0,2}\\
G\,\tfrac{M_S\,M_M}{r_\text{òrbita}^2} &= M_M\,\tfrac{v_\mathrm{orbital}^2}{r_\text{òrbita}} \Rightarrow v_\mathrm{orbital}^2 = \tfrac{G\,M_S}{r_\text{òrbita}} \quad\punts{0,2}
\end{aligned}\right\} \Rightarrow \]
\[ E_m = E_T = -\frac{1}{2}\,\frac{G\,M_S\,M_M}{r_\text{òrbita}} \quad\punts{0,2} = -3,90\cdot 10^{27}\ \mathrm{J} \quad\punts{0,2} \]
El signe negatiu significa que el satèl·lit Mimas està girant en una òrbita estable al voltant de Saturn. \punts{0,2}
\end{solapartat}
